\section{Tổng kết và hướng phát triển}

\subsection{Tổng kết kết quả nghiên cứu và hiện thực của đề tài}

Trải qua quá trình nghiên cứu lý thuyết, khảo sát nghiệp vụ thực tiễn, phân tích kiến trúc và hiện thực hóa phần mềm, đề tài \textbf{"Phát triển hệ thống quản lý văn phòng chia sẻ (CoSpace)"} đã hoàn thành xuất sắc toàn bộ các mục tiêu đặt ra cho đồ án tốt nghiệp ngành Khoa học Máy tính (CO4337 -- HK253). Hệ thống không chỉ dừng lại ở một công cụ quản lý cơ sở vật chất đơn thuần mà đã thực sự trở thành một nền tảng số toàn diện, hỗ trợ quản trị vận hành đa chi nhánh, tự động hóa dòng tiền và kiến tạo không gian kết nối chuyên môn cho cộng đồng thành viên.

So với giai đoạn Đồ án Chuyên ngành (ĐACN), hệ thống đã có những bước tiến mang tính đột phá về cả chiều sâu công nghệ lẫn tính hoàn thiện của sản phẩm:

\begin{enumerate}
    \item \textbf{Đột phá trong Tự động hóa Thanh toán (VietQR Napas 24/7):} Thay vì chỉ dừng lại ở các giao dịch mô phỏng ví điện tử, CoSpace đã tích hợp thành công cổng thanh toán Open Banking quốc gia (thông qua giải pháp PayOS). Hệ thống cung cấp mã VietQR động theo chuẩn Napas với cơ chế thẩm định chữ ký số HMAC-SHA256, tự động nhận diện giao dịch và cập nhật trạng thái đơn đặt chỗ sau 3-5 giây mà không cần người dùng tải lại trang.
    \item \textbf{Giải quyết triệt để bài toán tương tranh đụng lịch (Concurrency Control):} Xây dựng giải pháp phòng vệ hai tầng gồm Khóa giao dịch mức ứng dụng (\texttt{pg\_advisory\_xact\_lock}) và Ràng buộc loại trừ cấp cơ sở dữ liệu (\texttt{EXCLUDE USING gist}), ngăn chặn 100\% tình trạng đặt trùng chỗ ngồi (Double Booking) ngay cả khi có hàng trăm lượt truy cập đồng thời.
    \item \textbf{Ứng dụng Trí tuệ Nhân tạo thế hệ mới (Google Gemini AI):} Tích hợp mô hình ngôn ngữ lớn \texttt{gemini-3.5-flash} để vận hành Trợ lý ảo tư vấn thông minh (hỗ trợ giải đáp bảng giá, chính sách hoàn hủy 24/7), đồng thời kết hợp với thuật toán tương đồng Jaccard để tự động sinh lý do gợi ý kết nối đối tác mang tính cá nhân hóa cao.
    \item \textbf{Trải nghiệm người dùng trực quan và hiện đại:} Hiện thực Sơ đồ mặt bằng tương tác (SVG Interactive Floorplan Canvas) cho phép người dùng quan sát vị trí thực tế của bàn làm việc theo thời gian thực; Cung cấp bảng điều khiển Lễ tân hỗ trợ cơ chế dung sai 30 phút và quản lý dịch vụ phát sinh (Running Tab).
    \item \textbf{Đóng gói chuẩn mực và Triển khai Đám mây thực tế:} Áp dụng kỹ thuật Docker Multi-stage build rút gọn dung lượng Back-end xuống chỉ còn 215MB, triển khai thành công lên Render Cloud, phân phối Front-end trên mạng lưới biên Vercel Edge CDN và cơ sở dữ liệu trên Supabase Cloud.
    \item \textbf{Bảo chứng chất lượng với 250 ca kiểm thử tự động:} Hệ thống được bao phủ bởi 236 ca kiểm thử tự động Java (JUnit 5, Mockito) và 14 kịch bản kiểm thử tích hợp SQL chuyên sâu, đạt tỷ lệ thành công tuyệt đối 100\% PASS.
\end{enumerate}

\subsection{Đánh giá ưu điểm và mặt hạn chế của hệ thống}

\subsubsection{Ưu điểm nổi bật}
\begin{itemize}
    \item \textbf{Tính thực tiễn và nghiệp vụ chuyên sâu:} Hệ thống mô phỏng chuẩn xác toàn bộ các quy trình nghiệp vụ phức tạp của thị trường coworking space hiện đại: chính sách giá phân tầng linh hoạt (theo giờ/ngày/tháng), cơ chế hủy phạt 3 nấc thời gian minh bạch, bảo trì không gian định kỳ và đối soát tài chính tức thời.
    \item \textbf{Kiến trúc phần mềm chuẩn mực:} Phân tách rõ ràng giữa tầng hiển thị và tầng máy chủ (3-Tier Layered Architecture). Mã nguồn Back-end tuân thủ nghiêm ngặt các nguyên tắc Clean Architecture và SOLID, giúp mã nguồn dễ bảo trì, dễ mở rộng và có độ tương thích cao.
    \item \textbf{Tính an toàn và bảo mật cao:} Hệ thống triển khai cơ chế xác thực Stateless JWT kết hợp Supabase Auth, phân quyền truy cập nghiêm ngặt theo vai trò (RBAC), bảo vệ nguồn gốc chéo (CORS) và lưu vết kiểm toán (Audit Log) toàn vẹn mọi thao tác nhạy cảm.
\end{itemize}

\subsubsection{Những mặt còn hạn chế}
\begin{itemize}
    \item \textbf{Nền tảng người dùng:} Mặc dù giao diện web đã được tối ưu hóa theo hướng đáp ứng (Responsive Design) tương thích tốt trên điện thoại di động, hệ thống hiện vẫn chưa có ứng dụng di động bản địa (Native Mobile Application) chuyên biệt trên hai kho ứng dụng App Store và Google Play.
    \item \textbf{Tích hợp phần cứng thông minh:} Hiện tại quá trình check-in ra vào cửa vẫn phụ thuộc vào thao tác xác thực của nhân viên quầy lễ tân thông qua giao diện máy trạm, chưa tích hợp trực tiếp với các thiết bị phần cứng IoT (như khóa cửa từ tử thông minh, barrier cổng tự động).
\end{itemize}

\subsection{Hướng phát triển và mở rộng trong tương lai}

Để đưa hệ thống CoSpace tiến xa hơn nữa hướng tới mục tiêu thương mại hóa thực tế trên quy mô chuỗi dịch vụ, nhóm nghiên cứu đề xuất một số định hướng phát triển tiếp theo:

\begin{enumerate}
    \item \textbf{Phát triển Ứng dụng Di động Đa nền tảng (Mobile Native App):} Xây dựng ứng dụng di động bằng React Native hoặc Flutter nhằm tận dụng các tính năng phần cứng của điện thoại như Thông báo đẩy (Push Notifications) nhắc nhở lịch hẹn, quét mã QR check-in một chạm và mở khóa bằng sinh trắc học (FaceID / Fingerprint).
    \item \textbf{Tích hợp giải pháp Phần cứng Thông minh IoT (Smart IoT Access Control):} Kết nối máy chủ với các bộ điều khiển khóa cửa điện tử (Smart Door Locks) và cảm biến hiện diện tại vị trí làm việc thông qua giao thức mạng nhẹ \textbf{MQTT} hoặc \textbf{WebSockets}. Khi khách hàng thanh toán thành công, hệ thống tự động cấp mã số PIN tạm thời hoặc mã QR mở cửa có hiệu lực đúng trong khung giờ đã đặt.
    \item \textbf{Mở rộng Kiến trúc Đa bên thuê (Multi-tenant SaaS Architecture):} Nâng cấp kiến trúc dữ liệu và logic phân quyền để CoSpace trở thành một nền tảng phần mềm dưới dạng dịch vụ (SaaS). Cho phép các chủ doanh nghiệp văn phòng chia sẻ khác nhau có thể đăng ký tài khoản doanh nghiệp, tự thiết lập thương hiệu, cấu hình bảng giá và vận hành cơ sở kinh doanh của riêng mình trên cùng một hạ tầng máy chủ tập trung.
    \item \textbf{Tối ưu hóa gợi ý đối tác bằng Học máy Chuyên sâu (Advanced Graph AI):} Ứng dụng mạng nơ-ron đồ thị (Graph Neural Networks -- GNN) để mô hình hóa mạng lưới quan hệ giữa các thành viên, kết hợp dữ liệu lịch sử các dự án thành công để đưa ra các gợi ý hợp tác kinh doanh có độ chuẩn xác và giá trị thương mại cao hơn.
\end{enumerate}